\documentclass[12pt,a4paper]{article}
\usepackage{amsmath,amssymb,amsthm}
\usepackage{booktabs}
\usepackage{array}
\usepackage{geometry}
\usepackage{hyperref}
\usepackage{natbib}
\usepackage{setspace}
\usepackage{natbib}
\geometry{margin=2.5cm}
\onehalfspacing

\newtheorem{theorem}{Theorem}[section]
\newtheorem{proposition}[theorem]{Proposition}
\newtheorem{corollary}[theorem]{Corollary}
\newtheorem{lemma}[theorem]{Lemma}
\newtheorem{remark}[theorem]{Remark}
\theoremstyle{definition}
\newtheorem{definition}[theorem]{Definition}

\title{Loss Aversion, Endogenous Reference Points,\\
       and Boom-Bust Asymmetry in Financial Wealth Dynamics}
\author{Anurag Rathi\\
        \small Riskcare Ltd., London\\
        \small PhD Economics, University of Reading (2024)}
\date{April 2026 --- Working Paper}

\begin{document}
\maketitle

% ============================================================
% ABSTRACT
% ============================================================

\begin{abstract}
Financial markets exhibit a well-documented asymmetry between expansions and
contractions: boom periods are volatile, turbulent, and fast; bust periods are
compressed, slow, and grinding. This pattern is robust across asset classes,
time periods, and economies, yet lacks a structural explanation rooted in
individual preferences. This paper derives the asymmetry from first principles
using the minimal behavioural structure sufficient to deliver it as an
equilibrium theorem rather than a calibration choice.

We retain the single empirically robust element of prospect theory --- the kink
in the value function at a reference point --- and discard the rest. The
reference point is the economy-wide Schumpeterian wealth trend, publicly
observed and growing at a constant rate. These ingredients are embedded in a
Stackelberg market with a loss-averse market maker and a rational speculator,
scoped to liquid exchange-traded settings where adverse selection is attenuated
and dealer risk management is the primary source of price impact asymmetry.

Two structural features are jointly necessary. First, the zero-sum trading game
guarantees as a theorem that the two agents always sit on opposite sides of the
kink simultaneously, generating asymmetric price impact every period. Second,
the Schumpeterian trend provides mean-reversion that pins volatility persistence
to the economy's long-run growth rate. The paper derives a structural recursion
for the conditional variance of detrended wealth deviations, with persistence
determined by the growth rate alone and the asymmetry ratio by dealer loss
aversion alone --- a separation property that is the paper's primary
identification contribution. All structural mappings are derived analytically
and verified symbolically.
\end{abstract}


\noindent\textbf{Keywords:} loss aversion; boom-bust asymmetry; Threshold-GARCH; market microstructure; Stackelberg equilibrium; Schumpeterian growth

\newpage

% ============================================================
% SECTION 1: MOTIVATION
% ============================================================

\section{Motivation}
\label{sec:motivation}

\subsection{The empirical pattern this paper explains}
\label{sec:pattern}

One of the most robust patterns in financial history is the asymmetry between
boom and bust in the dynamics of financial wealth. Expansions --- periods when
aggregate dealer or investor wealth sits above its long-run growth trend ---
are characterised by rapid balance sheet growth, large position sizes, elevated
volatility, and turbulent price movements. Contractions --- periods of below-trend
wealth --- are slow, compressed, and grinding.

\citet{kindleberger1978} documents this pattern across two centuries of financial
manias and crashes: the run-up to a crisis is invariably turbulent and fast,
the aftermath slow and drawn out. \citet{adrian2010} show that financial intermediary balance sheets expand rapidly and with high variance when dealers
are above their funding benchmarks, and contract slowly when approaching their
limits. \citet{brunnermeier2009} establish that funding liquidity
constraints generate precisely this asymmetry at the position level: above
the constraint, dealer positions are large and wealth dynamics are turbulent;
at the constraint, positions compress and wealth process volatility falls
sharply. \citet{reinhart2009}, covering 66 countries across eight
centuries, document the same pattern at the macro level --- the run-up to a
financial crisis is associated with high and rising wealth volatility, while
the aftermath is characterised by persistent below-trend deviations with
compressed variance. \citet{danielsson2008} show theoretically and
empirically that endogenous risk --- generated by VaR-constrained intermediaries
responding to their own wealth positions --- is higher in expansions than
contractions, consistent with a structural variance asymmetry that has its
source in dealer behaviour rather than in the distribution of fundamentals.

The pattern is well described but structurally unexplained. Existing models
either take it as given --- describing it statistically without deriving it from
preferences --- or generate it from mechanisms that require strong distributional
assumptions about probability perception incompatible with rational market
behaviour. The present paper identifies the minimal structure that delivers the
asymmetry and its persistence as equilibrium theorems.

\subsection{Why Cumulative Prospect Theory is over-specified for this purpose}
\label{sec:cpt}

The natural behavioural candidate for a structural explanation is Cumulative
Prospect Theory (\citeauthor{tversky1992} (\citeyear{tversky1992})). CPT generates asymmetric
responses to gains and losses through two mechanisms: the kink in the value
function at the reference point, which makes losses loom larger than equivalent
gains, and the probability weighting function, which distorts small probabilities
upward and large probabilities downward. In static settings with exogenous
reference points, CPT is empirically well supported. In a dynamic financial
market equilibrium, however, it is analytically over-specified.

The probability weighting function is rank-dependent and non-differentiable,
making it incompatible with recursive dynamic modelling. Dynamic consistency
fails: a CPT agent's preferences over plans change as the plan is executed,
creating time-inconsistency that cannot be resolved without additional
assumptions. Aggregation across CPT agents with heterogeneous reference points
and probability weighting functions generates indeterminate equilibria.
Most damagingly for the present purpose, the reference point is either
exogenous --- which is indefensible in a market where agents form beliefs and
positions endogenously --- or endogenised through a rational expectations
fixed-point recursion (\citet{koszegi2006,koszegi2007}, 2006, 2007) that becomes
intractable when the reference point and equilibrium wealth volatility are
jointly determined.

\citet{ang2005} and \citet{routledge2010} show that many of
CPT's empirical successes can be replicated by retaining only the kink ---
through disappointment aversion or related specifications --- while discarding
the probability weighting function. This paper takes the same approach but
goes further: it also replaces the rational expectations reference point with
a simpler social benchmark, eliminating the fixed-point problem at its source.

The most closely related paper in the behavioural asset pricing literature
is \citet{barberis2001}, who place loss aversion over financial
wealth fluctuations at the centre of a representative agent consumption-based
asset pricing model and generate time-varying risk aversion through the house
money effect --- after prior gains the investor is less loss averse, after
prior losses more so. Four differences distinguish this paper from that
contribution. First, loss aversion here is fixed rather than time-varying:
the market maker's $\lambda$ does not change with prior performance. Second,
the reference point is a shared social Schumpeterian benchmark rather than
an individual historical benchmark that tracks personal investment outcomes.
Third, the model is a strategic market microstructure equilibrium with three
distinct agents rather than a representative agent consumption CAPM. Fourth,
the paper derives exact closed-form structural mappings verified symbolically
rather than characterising equilibrium properties through numerical simulation.
\citet{barberis2001} establish that loss aversion over financial
wealth can explain the equity premium and excess return volatility at the
aggregate level; this paper establishes what fixed dealer loss aversion
against a social benchmark delivers in a market microstructure equilibrium.

\subsection{The minimal structure}
\label{sec:minimal}

This paper focuses on a specific and empirically important class of markets:
liquid, exchange-traded settings where prices are publicly observable, adverse
selection is attenuated, and the dominant source of price impact asymmetry
is dealer risk management behaviour rather than private information. In these
markets, VaR constraints and balance sheet benchmarks govern dealer behaviour
more than informational disadvantage, and loss aversion is the structurally
appropriate primary mechanism.

We retain the single empirically robust element of prospect theory: the kink
in the value function at a reference point, which makes losses loom larger
than equivalent gains. Everything else --- the probability weighting function,
the rank-dependence, the S-shaped curvature away from the kink --- is discarded.
The reference point is the economy-wide Schumpeterian wealth trend, growing
at a constant rate $g > 0$ and publicly observed by all agents. This is the
simplest social benchmark consistent with a market setting. It requires no
individual wealth histories, no rational expectations fixed-point, and no
peer-group comparison. Its merit is that it is the weakest assumption capable
of delivering the structural results; its limitation is that it abstracts from
heterogeneous benchmarks and time-varying growth.

The theoretical motivation for a social reference point connects to the
social comparison literature \citep{festinger1954,fehr1999} and
to catching-up-with-the-Joneses models in asset pricing \citep{abel1990,campbell1999}. The Schumpeterian implementation is distinct from these
contributions: the reference point is endogenous to the system (determined
by the economy's growth trajectory) but exogenous to each individual agent's
decision, eliminating the strategic manipulation of reference points that
complicates \citet{koszegi2006,koszegi2007} implementations.

\subsection{Two necessary structural conditions}
\label{sec:necessary_conditions}

The structural results depend on two features of the minimal setup that are
jointly necessary: each is required for a specific structural result.

The first is the combination of the shared reference point with the zero-sum
trading structure and within-period position closing. Together, these three
conditions guarantee as a theorem --- Proposition~1 below --- that the market
maker and speculator always sit on opposite sides of the kink simultaneously.
When total detrended wealth is conserved, one agent's gain above the benchmark
is exactly the other's loss below it. The loss-aversion mechanism therefore
operates asymmetrically through the price-setter and the quantity-setter in
every period, generating the boom-bust variance asymmetry without additional
assumptions. If the shared reference point is replaced by own-history
benchmarks, this regime separation is lost and the structural link between
wealth volatility persistence and the growth rate disappears entirely.

The second is the Schumpeterian growth trend itself, which provides the
mean-reversion in the deviation recursion. Each period, detrended wealth
reverts toward the benchmark at rate $1/(1+g)$. This deflation factor is
the structural source of persistence in wealth volatility dynamics: it
connects the rate at which volatility shocks dissipate to the economy's
long-run growth rate in a precise and testable way. If the S-shaped value
function is removed --- that is, if the kink disappears at $\lambda = 1$ ---
the asymmetry between boom and bust variance vanishes exactly, confirming
that loss aversion is not only sufficient but necessary for the result.

The institutional counterpart of the model is a dealer trading desk operating
under Value-at-Risk constraints against an informed fund, both benchmarked
against the market's trend growth. \citet{li2025}, using
granular desk-level regulatory data on US Treasury dealers, document precisely
this mechanism: dealers reduce position size as they approach their internal
VaR limits, and accept lower prices to shed inventory when the limit binds.
This behaviour is isomorphic to the model's contraction regime, in which the
market maker --- above trend and protective of accumulated gains --- posts
high price impact and the speculator takes small positions.

\subsection{Contribution and relation to existing work}
\label{sec:contribution}

The paper makes four contributions. The first is the derivation of
boom-bust wealth volatility asymmetry as an equilibrium theorem from the
minimal CPT-like structure. The second is the structural identification
result: the persistence of wealth volatility is determined by the growth
rate alone, and the asymmetry ratio is determined by loss aversion alone,
with the two parameters orthogonal. This separation means that two observable
features of wealth volatility dynamics separately identify two economic
primitives without auxiliary assumptions. The third is the IGARCH boundary
result: as the growth rate approaches zero, volatility persistence approaches
one, providing the first preference-based microfoundation for integrated
volatility processes. The fourth is the connection of the structural mappings
to a testable cross-sectional prediction: in cross-country data, persistence
estimates should co-vary with long-run trend growth and asymmetry ratios
with measures of dealer loss aversion, with zero cross-effects.

The paper is postdictive in intent: it explains why the boom-bust asymmetry
in financial wealth takes the specific parametric form it does, given the
structural primitives of dealer loss aversion and trend growth, rather than
fitting or forecasting it. The natural empirical counterpart --- estimating
the structural recursion on detrended wealth or dealer balance sheet data
and testing whether the cross-sectional restrictions hold --- is a direction
for future work that the structural results are designed to motivate.

A fifth contribution concerns the relationship between loss aversion and
adverse selection. In markets where both channels operate simultaneously,
they are observationally entangled in the price impact coefficient: any
asymmetry in observed price impact could in principle reflect either the
market maker's informational disadvantage (adverse selection) or its
behavioural response to its own wealth position (loss aversion). By
shutting down adverse selection deliberately, this paper isolates the loss
aversion channel and establishes that it is sufficient for boom-bust
asymmetry even in the absence of private information. The separation
property --- $\partial\beta/\partial\lambda = 0$ and $\partial(|\gamma|/\alpha)/\partial g = 0$
--- then provides the identification strategy for a fuller model in which
both channels operate simultaneously: growth policy shifts $\beta$ without
affecting the asymmetry ratio, while macroprudential policy shifts
$|\gamma|/\alpha$ without affecting persistence. The two channels are
non-redundant and separately identified.

Existing micro-foundations for volatility clustering \citep{brock1998,follmer2005,watanabe2008} derive clustering
from heterogeneous agents or information structure rather than from individual
preferences in an explicit market equilibrium. The present paper is closest
in spirit to \citet{watanabe2008}, who embeds stochastic volatility in a
Kyle-style microstructure, but the volatility process there enters
exogenously through the information structure rather than being derived
from preferences.

\begin{remark}[Market quality across regimes]
\label{rem:market_quality}
\citet{kyle1985} characterises market quality through three dimensions:
depth (the inverse of price impact), tightness (the bid-ask spread), and
resilience (the speed of recovery after a shock). Kyle's model delivers a
single static value for each, determined by the signal-to-noise ratio
$\sigma_\mu/\sigma_u$. This paper extends all three into a regime-switching
setting driven by dealer loss aversion and the Schumpeterian growth rate.

Depth and tightness are regime-dependent. In expansion the effective price
impact is $\lambda^0_p/\lambda$ --- markets are deeper and tighter than
the Kyle benchmark by a factor of $\lambda$. In contraction it is
$\lambda \cdot \lambda^0_p$ --- shallower and wider by the same factor.
The ratio of expansion depth to contraction depth is $\lambda^2$,
providing a structural theory of liquidity procyclicality consistent with
\citet{adrian2010} and \citet{brunnermeier2009}.

Resilience is where the model adds most relative to Kyle, which does not
address it at all. The persistence parameter $\beta = 1/(1+g)^2$ is a
structural measure of recovery speed: high-growth economies have lower $\beta$
and faster recovery; stagnant economies approach $\beta = 1$, meaning
near-permanent impairment of market quality. This connects liquidity
resilience directly to macroeconomic growth --- a prediction invisible from
reduced-form microstructure estimation.
\end{remark}

% ============================================================
% SECTION 2: THE FRAMEWORK
% ============================================================

\section{The Framework}
\label{sec:framework}

\subsection{Players and roles}
\label{sec:players}

Three agents interact in a competitive financial market at each date $t$:

\begin{itemize}
  \item \textbf{Market maker (MM).} Posts the price impact schedule
    $\lambda_p > 0$ before trade. Absorbs total order flow. Has a
    loss-averse value function with kink parameter $\lambda > 1$ evaluated
    against the shared Schumpeterian benchmark. Acts as Stackelberg leader.

  \item \textbf{Speculator (S).} Receives private signal $\mu_t \sim
    \mathcal{N}(0,\sigma^2_\mu)$ about the asset's fundamental value and
    chooses position size $m_t \in \mathbb{R}$ after observing the posted
    price impact schedule. Loss-neutral --- standard expected profit
    maximiser. Acts as Stackelberg follower.

  \item \textbf{Noise traders.} The noise trader assumption follows
    \citet{kyle1985} and is standard throughout the market microstructure
    literature: without uninformed order flow, a competitive market-making
    equilibrium cannot exist because any order would perfectly reveal the
    speculator's signal and the market maker would always lose to the
    informed party. Noise traders submit random order flow
    $u_t \sim \mathcal{N}(0,\sigma^2_u)$, independent of $\mu_t$ and
    $\varepsilon_t$. The stationarity condition derived in
    Section~\ref{sec:calibration} requires the signal-to-noise ratio
    $\sigma_\mu/\sigma_u$ to exceed a threshold that depends on $\lambda$
    and $g$.
\end{itemize}

\begin{remark}[Uncertainty structure]
\label{rem:definetti}
The uncertainty structure of the model is de~Finetti subjective in character.
The private signal $\mu_t$ represents the speculator's epistemic advantage
--- knowledge the market maker does not possess --- rather than aleatory
randomness intrinsic to the world. The noise trader flow $u_t$ is closer to
aleatory uncertainty, representing the irreducible aggregation of many
independent liquidity-motivated trades whose individual motivations are
unobservable to all agents.

Crucially, the mechanism generating the boom-bust variance asymmetry is not
uncertainty-driven. Both agents observe the Schumpeterian trend and the
detrended wealth deviation $D_{t-1}$: the regime is public information.
The market maker's asymmetric price impact response is a certain reaction
to a publicly known position, not a probabilistic response to private
information. This distinguishes the paper's mechanism sharply from
\citet{kyle1985}, where asymmetric price impact arises from the market
maker's epistemic uncertainty about informed order flow. Here, loss aversion
operates on publicly observable states --- consistent with the predominance
of public price observation in exchange-traded markets where the information
asymmetry problem is attenuated and the risk management problem is primary.
\end{remark}

\begin{remark}[Why loss aversion must reside in the market maker]
\label{rem:mm_assignment}
A loss-averse speculator generates no equilibrium asymmetry. Under a
piecewise linear value function, the loss aversion parameter $\lambda$
appears as a positive scalar multiplying the speculator's objective and
cancels exactly in the first-order condition: the optimal position
$m^*_t = \mu_t/(2\lambda_p)$ is independent of $\lambda$. For loss
aversion to affect equilibrium outcomes it must enter through a binding
constraint rather than an interior optimum.

In the Stackelberg game, the market maker's loss aversion manifests as
self-imposed regime-dependent bounds on the price impact schedule it
posts. In the gain region (contraction regime, $D_{t-1} < 0$, MM above
the benchmark) the MM imposes an upper cap on price impact, protective
of accumulated wealth; in the loss region (expansion regime, $D_{t-1}
\geq 0$, MM below the benchmark) it accepts a lower floor, willing to
narrow margins to attract order flow and recover toward the benchmark.
Both bounds are derived from the MM's constrained optimisation and both
bind in equilibrium. Because the speculator's optimal position is the
signal divided by twice the posted price impact, any change in the price
impact schedule propagates directly and mechanically to the speculator's
position size. The leader's binding constraint propagates to the
follower's optimal choice; the follower's constraint does not propagate
to the leader. This is why asymmetry requires loss aversion in the
price-setter.

This assignment is independently grounded in the empirical microstructure
literature. \citet{li2025} document, using granular
desk-level regulatory data on US Treasury dealers, that dealers reduce
position changes and require higher compensation as they approach their
internal VaR limits --- precisely the high-price-impact, small-position
behaviour of the contraction regime. \citet{adrian2010,adrian2014}
document that dealer balance sheet dynamics are procyclical and
VaR-constrained at the aggregate level.
\end{remark}

\begin{remark}[Stackelberg timing]
\label{rem:stackelberg}
The Stackelberg ordering --- market maker as leader, speculator as
follower --- reflects the institutional reality that dealers post firm
quotes before customer order flow arrives, and those quotes constitute
binding obligations at the posted price for at least one trading increment
\citep{glosten1994,parlour2008}. The structural necessity of
this assignment is established in Remark~\ref{rem:mm_assignment}: assigning
loss aversion to the follower produces no asymmetry, so the Stackelberg
ordering is not merely institutionally motivated but structurally required.
\end{remark}

\subsection{The shared reference point}
\label{sec:reference}

The economy grows at a constant Schumpeterian rate $g > 0$. The shared
reference point is:
\begin{equation}
  \bar{r}_t = \bar{W}_0 \cdot (1+g)^t
  \label{eq:reference}
\end{equation}
where $\bar{W}_0 = (W^S_0 + W^{MM}_0)/2$ is the initial average wealth
of the two agents. This construction is publicly observed, continuously
evolving independently of individual positions, and exogenous to each
agent's decision while endogenous to the system. It is deliberately simple:
no individual wealth histories, no rational expectations fixed-point, and
no peer-group comparison.

The Schumpeterian growth trend serves two distinct roles. Its level
determines the threshold at which the regime switches --- the point at
which the MM moves from below-trend (expansion, loss region) to above-trend
(contraction, gain region). Its rate $g$ enters the deviation recursion as
the mean-reversion speed, pinning volatility persistence to the economy's
long-run growth trajectory.

\subsection{Detrended wealth and the state variable}
\label{sec:detrended}

Define detrended wealth for each agent by deflating nominal wealth by
the trend:
\begin{equation}
  \tilde{W}^S_t = \frac{W^S_t}{(1+g)^t}, \qquad
  \tilde{W}^{MM}_t = \frac{W^{MM}_t}{(1+g)^t}
\end{equation}
The state variable is the speculator's deviation from the shared benchmark:
\begin{equation}
  D_t \equiv \tilde{W}^S_t - \bar{W}_0
\end{equation}
The sign of $D_{t-1}$ determines the regime entering period $t$. When
$D_{t-1} \geq 0$ the speculator is above the benchmark (expansion regime)
and the MM is below it. When $D_{t-1} < 0$ the speculator is below the
benchmark (contraction regime) and the MM is above it. The zero-sum condition
--- that total detrended wealth $\tilde{W}^S_t + \tilde{W}^{MM}_t =
2\bar{W}_0$ is conserved at each date --- ensures that the two agents are
always in opposite regimes simultaneously. This is established formally in
Proposition~\ref{prop:regime_separation} and is the structural source of
the variance asymmetry.



\section{Price Process and Wealth Dynamics}
\label{sec:price}

Total order flow is $y_t = m_t + u_t$. Prices follow:
\begin{equation}
  p_t = p_{t-1} + \lambda_p y_t + \varepsilon_t
\end{equation}
where $\lambda_p > 0$ and $\mathbb{E}_{t-1}[\varepsilon_t] = 0$,
$\mathbb{E}_{t-1}[\varepsilon^2_t] = \sigma^2_t$. The \citet{kyle1985}
zero-profit benchmark pins:
\begin{equation}
  \lambda^0_p = \frac{\sigma_\mu}{2\sigma_u}
  \label{eq:kyle_lambda}
\end{equation}
The speculator's wealth under position closing is
$W^S_t = W^S_{t-1} + \lambda_p m^2_t + m_t\varepsilon_t$.
The zero-sum condition $W^S_t + W^{MM}_t = 2\bar{W}_0(1+g)^t$
gives $D_t + \tilde{\Delta}^{MM}_t = 0$ for all $t$.

\section{Reference Point, Regimes, and Regime Separation}
\label{sec:regimes}

Under the zero-sum condition, the two agents are always on opposite
sides of the reference point:
\begin{align}
  \text{Regime (+):} & \quad D_t \geq 0 \;\Leftrightarrow\;
    \tilde{\Delta}^{MM}_t \leq 0 \quad \text{(expansion)} \\
  \text{Regime (-):} & \quad D_t < 0 \;\Leftrightarrow\;
    \tilde{\Delta}^{MM}_t > 0 \quad \text{(contraction)}
\end{align}

\begin{proposition}[Regime Separation]
\label{prop:regime_separation}
Under zero-sum dynamics, position closing, and a shared reference point,
the two agents are never simultaneously in the loss region.
\end{proposition}

\begin{proof}
Under zero-sum dynamics $\tilde{\Delta}^S_t + \tilde{\Delta}^{MM}_t = 0$.
Simultaneous loss requires both deviations negative, implying their sum
is negative --- a contradiction.
\end{proof}

\section{Martingale Properties and Asymmetric Drift}
\label{sec:martingale}

\begin{proposition}[Price Martingale]
Under position closing and $\mathbb{E}_{t-1}[\varepsilon_t] = 0$, the
price process $\{p_t\}$ is a martingale.
\end{proposition}

\begin{proof}
From $p_t = p_{t-1} + \lambda_p m_t + \varepsilon_t$ and the MM's
competitive zero-profit condition, $\mathbb{E}_{t-1}[p_t] = p_{t-1}$.
\end{proof}


\begin{proposition}[Asymmetric Drift]
\label{prop:asym_drift_main}
Define $D^*_+ = \lambda\sigma^2_\mu(1+g)/(4g\lambda^0_p)$ and
$D^*_- = \sigma^2_\mu(1+g)/(4g\lambda\lambda^0_p)$, with $D^*_+/D^*_- = \lambda^2$.
Then: (i) in contraction, $\{D_t\}$ is a submartingale; (ii) in expansion,
$\{D_t\}$ is a submartingale for $D_{t-1} < D^*_+$ and supermartingale for
$D_{t-1} > D^*_+$; (iii) $D^*_+ > D^*_- > 0$: the economy has a positive
bias; (iv) the deflation factor $1/(1+g)$ is identical across regimes.
\end{proposition}

\begin{proof}
See Appendix, Proposition~\ref{prop:asym_drift}.
\end{proof}

\section{Counter-Cyclical Variance and the Boom-Bust Mechanism}
\label{sec:countercyclical}

The core mechanism is direct. In \textbf{expansion} ($D_{t-1} \geq 0$):
the MM is in the loss region and posts low price impact
$\lambda^+_p = \lambda^0_p/\lambda$; the speculator takes large positions
$m^{*+} = \lambda\mu_t/(2\lambda^0_p)$; the variance of $D_t$ is
proportional to $(B^+)^2 = \lambda^2/(4\lambda^{0\,2}_p)$.

In \textbf{contraction} ($D_{t-1} < 0$): the MM posts high price impact
$\lambda^-_p = \lambda\lambda^0_p$; the speculator takes small positions
$m^{*-} = \mu_t/(2\lambda\lambda^0_p)$; the variance of $D_t$ is
proportional to $(B^-)^2 = 1/(4\lambda^2\lambda^{0\,2}_p)$.

The variance ratio is $(B^+)^2/(B^-)^2 = \lambda^4$: expansions are
$\lambda^4$-times more volatile in $D_t$ than contractions. For
$\lambda = 2.25$ (Tversky-Kahneman), this is a factor of $25.6$.

This matches the boom-bust asymmetry documented by \citet{kindleberger1978}
and \citet{hamilton1989}: expansions are turbulent and fast, contractions are
compressed and grinding. Loss aversion is necessary and sufficient for
this asymmetry.

\section{Deviation Dynamics}
\label{sec:deviation}

From the wealth recursion and detrending:
\begin{equation}
  D_t = \frac{D_{t-1}}{1+g} + A^r\sigma^2_\mu + B^r\mu_t\varepsilon_t
  \label{eq:deviation_recursion}
\end{equation}
where in expansion $A^+ = \lambda/(4\lambda^0_p)$,
$B^+ = \lambda/(2\lambda^0_p)$, and in contraction
$A^- = 1/(4\lambda\lambda^0_p)$, $B^- = 1/(2\lambda\lambda^0_p)$.
The ratio $A^+/A^- = B^+/B^- = \lambda^2$.

Three channels operate: (i) Schumpeterian deflation $D_{t-1}/(1+g)$
--- the structural source of GARCH persistence $\beta$; (ii) positive
drift $A^r\sigma^2_\mu$ --- higher in expansion by factor $\lambda^2$;
(iii) stochastic component $B^r\mu_t\varepsilon_t$ --- higher variance
in expansion by factor $\lambda^4$.

\section{Equilibrium Derivation}
\label{sec:equilibrium}

\subsection{The MM's Problem and Regime-Dependent Price Impact}
\label{sec:mm_problem}

The MM maximises $\pi^{MM}(\lambda_p) = \lambda_p\sigma^2_u -
\sigma^2_\mu/(4\lambda_p)$ subject to regime-dependent risk limits.
The unconstrained optimum coincides with the Kyle zero-profit benchmark
$\lambda^0_p = \sigma_\mu/(2\sigma_u)$.

Loss aversion modifies the equilibrium through self-imposed binding
constraints: in the gain region (contraction), the MM imposes an upper
limit $\lambda_p \leq \lambda\lambda^0_p$; in the loss region
(expansion), the MM imposes a floor $\lambda_p \geq \lambda^0_p/\lambda$.
Both constraints bind, giving the regime-dependent price impact in equation~\eqref{eq:optimal_position}.

\subsection{Speculator's Optimal Position}
\label{sec:speculator}

The loss-neutral speculator solves $\max_{m_t} m_t\mu_t -
\lambda^{\rm eff}_p m^2_t$, giving:
\begin{equation}
  m^*_t = \frac{\mu_t}{2\lambda^{\rm eff}_p}
  = \begin{cases}
      \lambda\mu_t/(2\lambda^0_p) & D_{t-1} \geq 0 \\
      \mu_t/(2\lambda\lambda^0_p) & D_{t-1} < 0
    \end{cases}
  \label{eq:optimal_position}
\end{equation}
Position ratio: $m^{*+}/m^{*-} = \lambda^2$. The speculator's
FOC is identical in both regimes; all asymmetry flows from the MM's
pricing behaviour.

\subsection{The Cancellation Result}
\label{sec:cancellation}

The return innovation is $\varepsilon^r_t = \lambda^{\rm eff}_p m^*_t +
\varepsilon_t$. In each regime:
\begin{align*}
  \text{Expansion:} &\quad
    \frac{\lambda^0_p}{\lambda} \cdot \frac{\lambda\mu_t}{2\lambda^0_p}
    + \varepsilon_t = \frac{\mu_t}{2} + \varepsilon_t \\
  \text{Contraction:} &\quad
    \lambda\lambda^0_p \cdot \frac{\mu_t}{2\lambda\lambda^0_p}
    + \varepsilon_t = \frac{\mu_t}{2} + \varepsilon_t
\end{align*}
The signal component is $\mu_t/2$ in \emph{both} regimes. Therefore
$\operatorname{Var}(\varepsilon^r_t \mid \text{regime}) = \sigma^2_\mu/4 +
\sigma^2_\varepsilon \equiv \bar\sigma^2$ identically. The one-period
conditional variance of returns is regime-invariant. Variance structure in
$D_t$ emerges from the regime-dependent variance of the state variable
$(B^r)^2$, not from the conditional return variance.

\begin{remark}[No adverse selection baseline]
\label{rem:no_adv_sel}
The baseline model shuts down adverse selection deliberately. The return
innovation $\varepsilon^r_t = \mu_t/2 + \varepsilon_t$ is publicly
observable, and the cancellation result establishes that its conditional
variance is regime-invariant: the market maker faces identical
informational uncertainty in expansion and contraction. The variance
asymmetry therefore originates entirely from loss aversion --- the market
maker's behavioural response to its own publicly observable wealth
position relative to the Schumpeterian benchmark --- not from differential
information processing across regimes.

This is the appropriate baseline for the paper's identification purpose.
With both adverse selection and loss aversion present, the two channels
would be observationally entangled in the variance coefficients, making
clean separation of $\lambda$ and $g$ impossible without additional
structure. By isolating the loss aversion channel, the paper establishes
that it is sufficient for boom-bust asymmetry even in the absence of
private information. Robustness Choice~10 confirms that restoring adverse
selection scales $\alpha$ and $|\gamma|$ by a common factor $1/(1+\delta)$,
leaving $|\gamma|/\alpha = (\lambda^4-1)/\lambda^4$ intact: loss aversion
remains the source of asymmetry, adverse selection modulates its magnitude
without changing its structure.
\end{remark}

\begin{remark}[News interpretation]
\label{rem:news}
The cancellation result delivers a precise statement about news
interpretation in this model. The return innovation $\varepsilon^r_t =
\mu_t/2 + \varepsilon_t$ is the publicly observable news signal: its
variance is $\bar\sigma^2$ regardless of regime, so the same piece of
news arrives with identical statistical properties in expansion and
contraction. What differs across regimes is not the news itself but its
\emph{wealth consequence}. In expansion, the speculator's position is
$m^{*+} = \lambda\mu_t/(2\lambda^0_p)$; in contraction it is
$m^{*-} = \mu_t/(2\lambda\lambda^0_p)$. The ratio is $\lambda^2$: the
same signal $\mu_t$ generates a position --- and therefore a wealth
deviation --- that is $\lambda^2$ times larger in expansion than in
contraction. Loss aversion does not distort the interpretation of news;
it distorts the \emph{response} to news, and it does so asymmetrically
across the Schumpeterian threshold. This is the model's formalisation of
what news interpretation means in the minimal CPT-like structure.
\end{remark}

\subsection{The Threshold-GARCH Coefficients}
\label{sec:tgarch}


\begin{theorem}[Threshold-GARCH for Detrended Wealth Deviations]
\label{thm:main}
Under Assumptions (A1)--(A3), the conditional variance of $D_t$,
$V_t \equiv \operatorname{Var}(D_t \mid \mathcal{F}_t)$, satisfies:
\begin{equation}
  V_t = \omega + \alpha(\varepsilon^r_t)^2
       + \gamma(\varepsilon^r_t)^2\cdot\mathbf{1}[D_{t-1}<0]
       + \beta V_{t-1}
  \label{eq:tgarch}
\end{equation}
with exact structural identifications:
\begin{align}
  \alpha &= \Phi\lambda^2 > 0
    \label{eq:alpha} \\
  \gamma &= \Phi\,\frac{1-\lambda^4}{\lambda^2} < 0
    \quad \text{for all } \lambda > 1
    \label{eq:gamma} \\
  \beta  &= \frac{1}{(1+g)^2} \in (0,1)
    \quad \text{for } g > 0
    \label{eq:beta} \\
  \omega &= \frac{2\lambda^2\sigma^4_\varepsilon\sigma^4_\mu}
             {\lambda^{0\,2}_p(4\sigma^2_\varepsilon+\sigma^2_\mu)^2}
    \label{eq:omega}
\end{align}
where $\Phi \equiv \sigma^2_\varepsilon\sigma^2_\mu(16\sigma^4_\varepsilon
+\sigma^4_\mu) / [\lambda^{0\,2}_p(4\sigma^2_\varepsilon+\sigma^2_\mu)^3]$.

Key structural results: (i) $\gamma = 0$ iff $\lambda = 1$; (ii) $\gamma < 0$
for all $\lambda > 1$ --- expansion variance exceeds contraction variance;
(iii) $\beta \to 1$ as $g \to 0$ (IGARCH boundary); (iv) $|\gamma|/\alpha =
(\lambda^4-1)/\lambda^4$ depends only on $\lambda$; (v) $\beta$ depends only
on $g$.
\end{theorem}

\begin{proof}
The exact part follows from Proposition~\ref{prop:var_asym}: for each
realised regime $r$, $V^r_t = (B^r)^2\Sigma[H_0 + H_2(\varepsilon^r_t)^2]$.
Setting $\alpha = c_2^+$, $\gamma = c_2^- - c_2^+$, $\omega = c_0^+$, and
$\beta = 1/(1+g)^2$ from the Schumpeterian deflation in
\eqref{eq:deviation_recursion} gives \eqref{eq:tgarch}. Signs and the
$\beta$ limit follow from direct differentiation. Full derivation in
Proposition~\ref{prop:linearisation}.
\end{proof}

\begin{remark}[Sign of $\gamma$ and the boom-bust interpretation]
\label{rem:gamma_sign}
The asymmetry coefficient $\gamma = \Phi(1-\lambda^4)/\lambda^2$ is
strictly negative for all $\lambda > 1$. In the standard GJR-GARCH
literature, $\gamma > 0$ reflects the leverage effect: return variance is
higher after negative return shocks than after positive ones. The present
model has a different object: $V_t = \operatorname{Var}(D_t \mid
\mathcal{F}_t)$ is the conditional variance of \emph{detrended wealth
deviations}, not of returns. The correct benchmark is therefore not the
leverage effect but the boom-bust asymmetry in wealth dynamics.

With this benchmark, $\gamma < 0$ is the \emph{correct} prediction: boom
periods (expansion, $D_{t-1} \geq 0$) are more volatile than bust periods
(contraction, $D_{t-1} < 0$). This is precisely the pattern documented by
\citet{kindleberger1978} in his historical account of financial manias and by
\citet{hamilton1989} in his Markov-switching characterisation of business cycles:
expansions are turbulent and fast, contractions are slow and grinding. The
economic mechanism is transparent: in expansion, the loss-averse market
maker posts \emph{low} effective price impact ($\lambda^+_p =
\lambda^0_p/\lambda$), the speculator takes \emph{large} positions
($m^{*+} = \lambda\mu_t/(2\lambda^0_p)$), and the resulting variance
of $D_t$ is proportional to $(B^+)^2 = \lambda^2/(4\lambda^{0\,2}_p)$.
In contraction, the roles are reversed: high price impact, small positions,
variance proportional to $(B^-)^2 = 1/(4\lambda^2\lambda^{0\,2}_p)$. The
ratio is $(B^+)^2/(B^-)^2 = \lambda^4$.
\end{remark}

\label{sec:comparative_statics}

Theorem~1 delivers exact closed-form coefficients. The following
corollary characterises how those coefficients respond to the
underlying economic primitives $\lambda$, $g$, and $\lambda^0_p$.

\begin{corollary}[Comparative statics]
\label{cor:comp_statics}
The following partial derivatives hold exactly:
\begin{align}
  \frac{\partial\beta}{\partial g}
    &= -\frac{2}{(1+g)^3} < 0
    \label{eq:dbeta_dg} \\
  \frac{\partial|\gamma|}{\partial\lambda}
    &= \frac{2\Phi(\lambda^4+1)}{\lambda^3} > 0
    \label{eq:dgamma_dlam} \\
  \frac{\partial|\gamma|}{\partial g}
    &= 0
    \label{eq:dgamma_dg} \\
  \frac{\partial\alpha}{\partial\lambda}
    &= \frac{2\Phi\lambda} > 0
    \label{eq:dalpha_dlam} \\
  \frac{\partial\alpha}{\partial g}
    &= 0
    \label{eq:dalpha_dg} \\
  \frac{\partial^2\beta}{\partial g^2}
    &= \frac{6}{(1+g)^4} > 0
    \label{eq:d2beta_dg2}
\end{align}
Moreover, $\partial\alpha/\partial\lambda^0_p < 0$ and
$\partial|\gamma|/\partial\lambda^0_p < 0$: both shock-sensitivity
coefficients are decreasing in market liquidity.
\end{corollary}

\begin{proof}
All derivatives follow by direct differentiation of the
closed-form expressions in Theorem~1, using $\partial\Phi/\partial g
= 0$ (since $\Phi$ contains only $\sigma_\mu$, $\sigma_\varepsilon$,
$\lambda^0_p$, and $\lambda$). Each sign is immediate from the
positivity of all parameters and the facts that $\lambda > 1$ and
$g > 0$. All expressions are verified in the companion SymPy notebook.
\qed
\end{proof}

Four economic implications follow directly from Corollary
\ref{cor:comp_statics}.

\paragraph{Separation property.}
The most striking implication is a clean orthogonality: $\beta$
depends on $g$ and $g$ alone; $|\gamma|/\alpha =
(\lambda^4-1)/\lambda^4$ depends on $\lambda$ and $\lambda$ alone:
\[
\frac{\partial\beta}{\partial\lambda} = 0, \qquad
\frac{\partial(|\gamma|/\alpha)}{\partial g} = 0.
\]
Both are verified in the companion notebook. The persistence
parameter $\beta$ is therefore a sufficient statistic for the
Schumpeterian growth rate $g$, and the asymmetry ratio
$|\gamma|/\alpha$ is a sufficient statistic for the loss aversion
parameter $\lambda$. The two deep parameters are \emph{separately
identified} from standard GJR-GARCH output, with no need for
auxiliary structural assumptions beyond the model. This clean
separation is the paper's primary identification contribution.

\paragraph{Orthogonal policy channels.}
Equations \eqref{eq:dbeta_dg}--\eqref{eq:dgamma_dg} imply that
growth-enhancing policies (raising $g$) reduce volatility
persistence $\beta$ without affecting the asymmetry $|\gamma|$;
dealer risk-constraint policies (changing the effective $\lambda$
via VaR limits or capital requirements) alter $|\gamma|$ without
affecting $\beta$. The two channels are orthogonal, making
identification from observational data cleaner than in models where
all GARCH parameters respond to every primitive simultaneously.

\paragraph{Convexity of the persistence effect.}
The second derivative \eqref{eq:d2beta_dg2} is positive, meaning
the impact of growth on persistence accelerates as $g$ approaches
zero. The marginal benefit of growth for volatility stability is
largest in stagnant economies --- a structural argument for
prioritising growth policy in low-growth environments.

\paragraph{Liquidity as a common scaling factor.}
Both $\alpha$ and $|\gamma|$ are decreasing in $\lambda^0_p$, but
their ratio $|\gamma|/\alpha = (\lambda^4-1)/\lambda^4$ is
independent of $\lambda^0_p$. Market liquidity scales the level of
all variance coefficients without affecting their relative magnitude.
This means the asymmetry ratio $|\gamma|/\alpha$ can be used to
identify $\lambda$ from cross-market data without requiring
knowledge of or controls for market liquidity --- a useful robustness
property for empirical work.

Table~\ref{tab:comp_statics} presents numerical comparative statics
at the baseline parameterisation.

\begin{table}[ht]
\centering
\caption{Numerical comparative statics. Base:
$\sigma_\mu=1$, $\sigma_\varepsilon=0.5$, $\lambda^0_p=0.8$,
$g=0.02$.}
\label{tab:comp_statics}
\begin{tabular}{r cccc}
\toprule
& $\alpha$ & $|\gamma|$ & $\beta$ & $|\gamma|/\alpha$ \\
\midrule
\multicolumn{5}{l}{\textit{Varying $\lambda$ ($g=0.02$):}} \\
$\lambda=1.0$ & 0.0977 & 0.0000 & 0.9612 & 0.000 \\
$\lambda=1.5$ & 0.2197 & 0.1763 & 0.9612 & 0.803 \\
$\lambda=2.0$ & 0.3906 & 0.3662 & 0.9612 & 0.938 \\
$\lambda=2.25$& 0.4944 & 0.4751 & 0.9612 & 0.961 \\
$\lambda=3.0$ & 0.8789 & 0.8681 & 0.9612 & 0.988 \\
\midrule
\multicolumn{5}{l}{\textit{Varying $g$ ($\lambda=2.25$):}} \\
$g=0\%$  & 0.4944 & 0.4751 & 1.0000 & 0.961 \\
$g=1\%$  & 0.4944 & 0.4751 & 0.9803 & 0.961 \\
$g=2\%$  & 0.4944 & 0.4751 & 0.9612 & 0.961 \\
$g=5\%$  & 0.4944 & 0.4751 & 0.9070 & 0.961 \\
$g=10\%$ & 0.4944 & 0.4751 & 0.8264 & 0.961 \\
\bottomrule
\end{tabular}
\medskip

\noindent\footnotesize{\textit{Note:} $\alpha$ and $|\gamma|$ are
invariant to $g$; $\beta$ and $|\gamma|/\alpha$ are invariant to
$\lambda$. The separation property holds exactly.}
\end{table}

\section{Calibration and the Stationarity Requirement}
\label{sec:calibration}

Proposition~\ref{prop:stationarity} establishes that the Threshold-GARCH
process for $V_t$ is covariance-stationary if and only if $\Phi <
\Phi_{\max}$, where
\begin{equation}
  \Phi_{\max} \;=\; \frac{2g\lambda^2(g+2)}{(1+g)^2(\lambda^4+1)}.
  \label{eq:Phi_max}
\end{equation}
Substituting the Kyle benchmark $\lambda^0_p = \sigma_\mu/(2\sigma_u)$,
the condition translates into a signal-to-noise requirement on the ratio
$\sigma_\mu/\sigma_u$ of informed signal volatility to noise trader flow
volatility. The stationarity condition is therefore a \emph{liquidity
condition}: markets with insufficient informed trading relative to noise
trading, combined with high dealer loss aversion and low growth, produce
non-stationary volatility in detrended wealth deviations.

Table~\ref{tab:calibration} reports the minimum $\sigma_\mu/\sigma_u$
required for stationarity across a grid of $(\lambda, g)$ values, computed
at $\sigma_\mu = 1$ and $\sigma_\varepsilon = 0.5$.

\begin{table}[ht]
\centering
\caption{Minimum $\sigma_\mu/\sigma_u$ required for covariance-stationarity
         of $V_t$. Computed at $\sigma_\mu=1$, $\sigma_\varepsilon=0.5$.}
\label{tab:calibration}
\begin{tabular}{r ccccc}
\toprule
$\lambda$ & $g=1\%$ & $g=2\%$ & $g=3\%$ & $g=5\%$ & $g=10\%$ \\
\midrule
1.50 & 4.13 & 2.95 & 2.42 & 1.90 & 1.39 \\
2.00 & 5.19 & 3.70 & 3.04 & 2.39 & 1.75 \\
2.25 & 5.78 & 4.11 & 3.38 & 2.66 & 1.95 \\
2.50 & 6.38 & 4.54 & 3.74 & 2.94 & 2.15 \\
3.00 & 7.60 & 5.42 & 4.45 & 3.50 & 2.56 \\
\bottomrule
\end{tabular}
\end{table}

At the Tversky-Kahneman parameter value $\lambda = 2.25$ and a standard
annual equity growth rate of $g = 2\%$, stationarity requires
$\sigma_\mu/\sigma_u > 4.11$. Empirical estimates of the permanent price impact of trades --- the
empirical counterpart of the model's signal-to-noise ratio
$\sigma_\mu/\sigma_u$ in the Kyle tradition --- consistently exceed this
threshold. \citet{hasbrouck1991} reports estimates of the information
content of trades in the range 3--20 for NYSE stocks at the trade-by-trade
level, where the information content measures the permanent component of
price impact attributable to informed order flow rather than the adverse
selection spread component. Large-capitalisation equities, which are the natural empirical counterpart
for the model, lie comfortably within the stationary region. The model is
therefore empirically operational at the parameter values most relevant for
its intended application. The Schumpeterian trend $\bar{r}_t$ plays the role
of \citet{hasbrouck1991}'s quote midpoint in this framework: both are the
publicly observable efficient benchmark from which transaction-level
deviations are measured and to which they mean-revert. The mean-reversion
speed $1 - \beta = 1 - 1/(1+g)^2$ in the deviation recursion corresponds
directly to Hasbrouck's transient price impact component.

\begin{remark}[VAR structure]
\label{rem:var}
The deviation recursion $D_t = D_{t-1}/(1+g) + A^r\sigma^2_\mu +
B^r \mu_t\varepsilon_t$ is structurally an AR(1) with autoregressive
coefficient $1/(1+g)$ and regime-dependent innovation. This is the
structural counterpart of the bivariate VAR that \citet{hasbrouck1991}
estimates empirically: the permanent component of any shock to $D_t$
is zero under stationarity ($g > 0$) because all wealth deviations
eventually revert to the Schumpeterian benchmark, while the transient
component decays geometrically at rate $1/(1+g)$ per period. At the
IGARCH boundary ($g \to 0$) the autoregressive coefficient approaches
one and shocks become permanent, giving a unit-root VAR. The Threshold-GARCH
variance recursion for $V_t$ then overlays a regime-dependent second-moment
structure on top of this AR foundation --- an extension beyond Hasbrouck's
constant-variance framework. The natural empirical strategy is therefore a
regime-switching extension of the \citet{hasbrouck1991} VAR with threshold
at the Schumpeterian trend crossing $D_{t-1} = 0$, estimating $\beta$ from
the autoregressive coefficient and the variance ratio $\lambda^4$ from the
regime-dependent innovation variances.
\end{remark}

Three further properties of \eqref{eq:Phi_max} are worth noting. First,
$\Phi_{\max} \to 0$ as $g \to 0$: for any fixed $\Phi > 0$, the
stationarity condition eventually fails as growth vanishes. This is precisely
the IGARCH boundary of Proposition~\ref{prop:igarch} --- non-stationary
volatility in detrended wealth deviations is the structural signature of a
stagnant Schumpeterian economy. Second, $\Phi_{\max}$ is decreasing in
$\lambda$: more loss-averse dealers impose a higher liquidity requirement
for stationarity, reflecting the amplified regime asymmetry encoded in the
$\lambda^4$ ratio of Proposition~\ref{prop:Vt_ratio}. Third, at $\lambda = 1$
(loss-neutral dealers), $\Phi_{\max} = g(g+2)/(1+g)^2$, which is positive
for all $g > 0$ and easily achievable, confirming that the stationarity
constraint is entirely a consequence of dealer loss aversion rather than a
feature of the Schumpeterian growth structure alone.

The persistence parameter $\beta = 1/(1+g)^2$ is plotted against $g$ in
the text and takes the values $0.980$, $0.961$, $0.943$, $0.907$, and
$0.826$ at $g \in \{1\%, 2\%, 3\%, 5\%, 10\%\}$ respectively. These are
consistent with empirical GARCH estimates for equity markets, which
typically find $\beta \in (0.85, 0.98)$, and provide a structural
interpretation: markets with higher long-run growth rates should exhibit
lower volatility persistence, a prediction that is testable in cross-country
panel data.

% ============================================================
% SECTION 11: ROBUSTNESS
% ============================================================

\section{Robustness: Verifying the Necessary Structural Conditions}
\label{sec:robustness}

The structural results rest on two necessary conditions: the shared
Schumpeterian reference point (which, together with zero-sum dynamics and
position closing, is necessary for regime separation) and loss aversion
$\lambda > 1$ in the value function (which is necessary and sufficient
for the variance asymmetry). This section examines how the boom-bust
asymmetry and the persistence identification are affected when each of
the eleven design choices is relaxed, verifying that these two conditions
are necessary rather than merely convenient. The key finding is that
$\partial|\gamma|/\partial\lambda > 0$ --- more dealer loss aversion
always produces more boom-bust asymmetry --- survives all eleven
relaxations. The $\partial\beta/\partial g < 0$ direction fails only
when the two necessary conditions are directly removed.

\subsection*{Effect of each design choice on the structural results}

\noindent\textbf{Choice 1 --- Linear price impact.} Nonlinear price
impact breaks the exact cancellation result: the product
$\lambda^{\rm eff}_p \cdot m^*_t$ is no longer $\mu_t/2$ in both regimes.
The variance recursion becomes a Power ARCH (PARCH) specification.
The boom-bust variance asymmetry survives because the regime-dependent
bound structure is preserved; both comparative static directions are
maintained.

\noindent\textbf{Choice 2 --- Quadratic cost (price impact).} Without
price impact ($\lambda^0_p \to 0$), the speculator's problem has no
finite interior solution and the $\beta$ channel collapses. Result:
Threshold-ARCH(1). The Schumpeterian mean-reversion channel --- the
structural source of $\beta = f(g)$ --- is disrupted. The
$\partial|\gamma|/\partial\lambda > 0$ direction is preserved. This
choice is \emph{a necessary condition for the persistence identification}.

\noindent\textbf{Choice 3 --- Position closing.} Carry-over positions
introduce additional lags, extending the recursion to Threshold-GARCH
$(n,1)$. Both comparative static directions are preserved in each lag.

\noindent\textbf{Choice 4 --- Zero-sum condition.} Spread income $s > 0$
breaks zero-sum and admits a third regime in which both agents can
simultaneously be on the same side of the kink. Result: three-regime
asymmetric variance recursion. The $\partial\beta/\partial g < 0$
direction is preserved; the loss aversion effect is enriched with three
separate $\lambda$-dependent asymmetry terms.

\noindent\textbf{Choice 5 --- Regime structure.} The regime structure is
inherited from Choices 3, 4, and 6. When the shared reference point is
intact, regime separation holds and the variance asymmetry is preserved
regardless of the number of lags or the presence of spread income.

\noindent\textbf{Choice 6 --- Shared reference point.} Replacing the
shared Schumpeterian trend with own-history reference points eliminates
the Schumpeterian deflation from the deviation recursion. The
persistence parameter becomes $\beta = \rho^2$ where $\rho$ is the
own-history mean-reversion coefficient, severing the structural link
between $\beta$ and $g$. This choice is \emph{a necessary condition for the
persistence identification}: removing it eliminates the paper's second
main result.

\noindent\textbf{Choice 7 --- Identical $\lambda$.} Heterogeneous loss
aversion $\lambda^{MM} \neq \lambda^S$: the variance asymmetry is driven
by $\lambda^{MM}$ alone, since $\lambda^S$ cancels in the speculator's
first-order condition (Remark~\ref{rem:mm_assignment}). Both comparative
static directions are preserved.

\noindent\textbf{Choice 8 --- Constant $g$.} Time-varying $g_t$
delivers $\beta_t = 1/(1+g_t)^2$ period by period, yielding a
Markov-Switching Threshold-GARCH with growth cycles as the state variable.
The separation property holds period by period: $\beta_t$ depends only
on $g_t$ and the asymmetry ratio depends only on $\lambda$.

\noindent\textbf{Choice 9 --- MM zero-profit benchmark.} A monopolist MM
yields an identical $\lambda^0_p$ in the linear Kyle framework. Both
comparative static directions are fully preserved.

\noindent\textbf{Choice 10 --- Public signal.} Under private information,
adverse selection scales both $\alpha$ and $|\gamma|$ uniformly by
$1/(1+\delta)$ where $\delta$ is the adverse selection premium. The
ratio $|\gamma|/\alpha = (\lambda^4-1)/\lambda^4$ is preserved
exactly, as is the separation property.

\noindent\textbf{Choice 11 --- Stationarity ($g > 0$).} As $g \to 0$,
$\beta \to 1$ (IGARCH boundary, Proposition~\ref{prop:igarch}). The
$\partial\beta/\partial g$ derivative diverges at $g = 0$, reflecting
that the persistence effect is strongest precisely in stagnant economies.
The $\partial|\gamma|/\partial\lambda > 0$ direction is preserved, and
$|\gamma|$ is independent of $g$ (verified symbolically).

\medskip\noindent
Table~\ref{tab:robustness_expanded} summarises the robustness of both
comparative static directions across all eleven choices.

\begin{table}[ht]
\centering
\caption{Robustness of design choices: effect on structural identifications.}
\label{tab:robustness_expanded}
\small
\begin{tabular}{p{0.8cm} p{2.8cm} p{2.4cm} p{2.4cm} p{2.3cm}}
\toprule
Choice & Relaxation & Variance recursion
& $\partial\beta/\partial g < 0$?
& $\partial|\gamma|/\partial\lambda > 0$? \\
\midrule
1 & Nonlinear price impact   & PARCH          & Preserved & Preserved \\
2 & No quadratic cost        & Threshold-ARCH & \textbf{Disrupted} & Preserved \\
3 & Carry-over positions     & T-GARCH$(n,1)$ & Preserved & Preserved \\
4 & Spread income ($s>0$)    & 3-regime GARCH & Preserved & Enriched \\
5 & Regime structure         & Inherited      & Inherited & Inherited \\
6 & Own-history ref point    & $\beta=\rho^2$ & \textbf{Disrupted} & Preserved \\
7 & Het.\ $\lambda$          & T-GARCH        & Preserved & Preserved \\
8 & Time-varying $g_t$       & MS-T-GARCH     & Preserved (p-b-p) & Preserved \\
9 & Monopolist MM            & T-GARCH        & Preserved & Preserved \\
10 & Private signal          & Scaled T-GARCH & Preserved & Preserved \\
11 & $g \to 0$               & T-IGARCH       & Boundary  & Preserved \\
\bottomrule
\end{tabular}
\medskip

\noindent\footnotesize{\textit{Note:} T-GARCH = Threshold-GARCH; MS = Markov-Switching;
p-b-p = period by period. Choices~2 and~6 remove the necessary conditions for
the $\beta = f(g)$ identification. The $\partial|\gamma|/\partial\lambda > 0$
direction --- more dealer loss aversion produces more boom-bust asymmetry ---
survives all eleven relaxations, confirming it is the most robust result in
the paper. Both conditions are necessary rather than merely convenient:
removing either eliminates one of the two main structural identifications.}
\end{table}

% ============================================================
% SECTION 12: CONCLUSION
% ============================================================

\section{Conclusion}
\label{sec:conclusion}

\subsection{Summary}
\label{sec:summary}

This paper identifies the minimal behavioural structure sufficient to
derive the boom-bust asymmetry in financial wealth dynamics as an
equilibrium theorem. We retain the single empirically robust element of
prospect theory --- the kink in the value function at a reference point
--- and embed it in a Stackelberg market with a Schumpeterian background
growth trend as the shared reference point.

Two structural features are jointly necessary conditions. First, the combination
of zero-sum dynamics, the shared reference point, and within-period
position closing guarantees that the market maker and speculator always
sit on opposite sides of the kink simultaneously --- a derived theorem,
not an assumption. This generates the boom-bust variance asymmetry:
expansions are $\lambda^4$ times more volatile in detrended wealth
deviations than contractions, where $\lambda$ is the dealer's loss
aversion parameter. Second, the Schumpeterian trend provides the
mean-reversion that pins volatility persistence to the economy's long-run
growth rate $g$ through the exact structural mapping $\beta = 1/(1+g)^2$.

The separation property is the paper's primary identification result:
$\beta$ depends on $g$ alone (independent of $\lambda$) and
$|\gamma|/\alpha = (\lambda^4-1)/\lambda^4$ depends on $\lambda$ alone
(independent of $g$). The two structural primitives --- growth and dealer
loss aversion --- are separately identified from the shape and persistence
of observed wealth volatility dynamics, without auxiliary assumptions.

The paper began with the goal of finding micro-foundations of GARCH
behaviour through a minimal CPT-like structure. That goal is achieved, but
with a precise qualification: the GARCH-type recursion is derived for
the conditional variance of detrended wealth deviations rather than
directly for equity return variance, which is the conventional object of
empirical GARCH estimation. This is not a failure of the approach but a
clarification of what the minimal structure delivers. The cancellation
result --- that return variance is regime-invariant --- shows that the
same piece of news arrives identically in both regimes; what differs is
the wealth consequence of that news, which is $\lambda^2$ times larger in
expansion than in contraction. Boom-bust asymmetry in financial wealth is
therefore the correct and precise object of the micro-foundation, and the
paper delivers it as an equilibrium theorem rather than a calibration
choice. The connection to return variance is a direction for future work.

\subsection{The IGARCH boundary and its interpretation}
\label{sec:igarch}

As $g \to 0$, $\beta \to 1$ and the conditional variance of detrended
wealth deviations follows a Threshold-IGARCH process. This provides the
first preference-based micro-foundation for integrated volatility
processes. Near-unit-root volatility persistence --- documented in
low-growth episodes such as post-2008 Japan and post-crisis Europe ---
is not an arbitrary statistical restriction but the volatility regime of
a stagnant Schumpeterian economy.

Importantly, the asymmetry coefficient $\gamma$ is independent of $g$:
$\partial\gamma/\partial g = 0$. A stagnant economy therefore exhibits
the same boom-bust asymmetry as a growing one, but with permanently
persistent volatility shocks. The two features --- asymmetry and
persistence --- are independently determined by independent structural
primitives.

\subsection{Why loss aversion must reside in the market maker}
\label{sec:why_mm}

The assignment of loss aversion to the market maker is not a modelling
choice but a structural necessity established in
Remark~\ref{rem:mm_assignment}: loss aversion in the speculator's
objective cancels in the first-order condition and generates no
asymmetry. The mechanism requires loss aversion to enter through a
binding constraint on the price-setter, which then propagates to the
follower's optimal quantity choice.

This is independently confirmed at the desk level by \citet{li2025}, who document that US Treasury dealers reduce position
changes and require higher compensation as they approach their internal
VaR limits. The model's contraction regime --- high price impact, small
speculator positions, compressed wealth volatility --- matches this
documented behaviour precisely.

\subsection{Testable predictions and the identification programme}
\label{sec:predictions}

\noindent\textbf{Prediction 0} (empirical strategy). The structural
recursion can be taken to data using a regime-switching extension of the
\citet{hasbrouck1991} VAR framework. Estimate a bivariate system in
detrended wealth or dealer balance sheet deviations and regime indicator,
with threshold at the Schumpeterian trend crossing $D_{t-1} = 0$. The
autoregressive coefficient in the deviation equation identifies
$\beta = 1/(1+g)^2$ and the ratio of regime-specific innovation variances
identifies $\lambda^4$. The cross-sectional restriction --- $\beta$
co-varies with long-run TFP growth and $\lambda^4$ co-varies with
experimental loss aversion estimates across countries, with zero
cross-effects --- is the primary test of the separation property.

The separation property generates a falsifiable cross-sectional
restriction invisible from single-market time-series estimation. Across
markets $i = 1, \ldots, N$, the model predicts:
\begin{equation}
  \hat{\beta}_i = \frac{1}{(1+g_i)^2} + u_i, \qquad
  \frac{|\hat{\gamma}_i|}{\hat{\alpha}_i}
    = \frac{\lambda_i^4 - 1}{\lambda_i^4} + v_i
  \label{eq:cross_sectional}
\end{equation}
where the two equations share no common primitive: $g_i$ enters only
the first and $\lambda_i$ only the second. The natural empirical test
uses Penn World Tables long-run TFP growth for $g_i$, cross-cultural
experimental estimates of loss aversion for $\lambda_i$ (Wang et al.,
2017, provide 53-country evidence), and Threshold-GARCH estimation for
$\hat{\beta}_i$, $\hat{\alpha}_i$, $|\hat{\gamma}_i|$ on detrended
wealth or dealer balance sheet data.

The policy implication of the orthogonality is direct. Growth-enhancing
structural reforms reduce wealth volatility persistence $\beta$ without
affecting asymmetry $|\gamma|/\alpha$. Macroprudential policies that
constrain dealer risk-taking --- tighter VaR limits, higher capital
requirements --- reduce boom-bust asymmetry without affecting persistence.
The two policy instruments operate through different structural channels
and are therefore non-redundant: a result that cannot be derived from
reduced-form estimation.

\subsection{Policy implications}
\label{sec:policy}

The orthogonality of the two structural identifications generates a
precise policy implication. Because $\beta = f(g)$ only and
$|\gamma|/\alpha = f(\lambda)$ only, growth-enhancing structural reforms
and macroprudential policies that constrain dealer risk-taking are
\emph{non-redundant instruments}: they target different features of wealth
volatility dynamics through different structural channels.

This connects to the practitioner debate on procyclicality in financial
markets. \citet{adrian2010,adrian2014} document that
dealer balance sheets expand rapidly above trend and contract slowly
below it --- precisely the $\lambda^4$ variance asymmetry the model
quantifies. The model says that the amplitude of this procyclicality
is governed by dealer loss aversion $\lambda$, not by the growth rate.
A macroprudential policy that reduces effective $\lambda$ --- through
tighter VaR limits, higher capital requirements, or countercyclical
capital buffers --- reduces $|\gamma|/\alpha$ without affecting $\beta$.
Conversely, a growth-enhancing reform that raises $g$ reduces $\beta$
--- shortening the half-life of volatility shocks in wealth dynamics ---
without affecting the boom-bust asymmetry.

The structural model therefore provides a micro-founded rationale for
why both instruments are simultaneously necessary: macroprudential
policy addresses the cross-sectional asymmetry ($\lambda$ channel) while
growth policy addresses persistence ($g$ channel). Neither instrument
substitutes for the other, and the model quantifies the magnitude of
each effect through the exact structural mappings of Theorem~1 and
Corollary~\ref{cor:comp_statics}. This is a structural claim about
the architecture of financial volatility that reduced-form estimation
cannot deliver.

\subsection{Limitations and future directions}
\label{sec:limitations}

Three limitations should be stated clearly.

First, the persistence identification is a leading-order approximation.
The exact expression for $\beta$ involves the normal density evaluated
at the contraction threshold, and the approximation error is bounded at
under 6\% for moderate deviations from the stationary point (Proposition
A.5). The structural mapping $\beta = 1/(1+g)^2$ is the leading-order
term; all qualitative results and the separation property hold exactly.

Second, the paper derives a structural recursion for the conditional
variance of \emph{detrended wealth deviations} $V_t =
\operatorname{Var}(D_t \mid \mathcal{F}_t)$, not for the conditional
variance of equity returns. The cancellation result (Section
\ref{sec:cancellation}) establishes that the one-period return variance
is $\bar\sigma^2$ in both regimes --- regime-invariant --- so the
structural asymmetry operates through wealth dynamics rather than return
dynamics. The connection between wealth deviation variance and return
variance, while structurally motivated, requires a bridge argument that
is left for future work. The object of the paper --- boom-bust asymmetry
in financial wealth --- is the direct empirical counterpart of the
documented patterns in dealer balance sheet dynamics (Adrian and Shin,
2010, 2014; Brunnermeier and Pedersen, 2009) and in post-crisis wealth
recovery paths (Reinhart and Rogoff, 2009).

Third, the reference point construction --- a common Schumpeterian trend
at a constant growth rate --- abstracts from heterogeneous benchmarks,
time-varying growth, and catching-up dynamics. Time-varying $g_t$
following a Schumpeterian wave process is the most important extension:
it delivers Markov-Switching Threshold-GARCH with growth cycles as the
state variable, connecting the framework to Hamilton-style regime
switching in growth rates. Heterogeneous loss aversion $\lambda^{MM}
\neq \lambda^S$ and carry-over positions (extending to Threshold-GARCH
$(n,1)$) are further natural extensions.


\appendix

\section{Technical Appendix}
\label{app:proofs}

This appendix collects all proofs omitted from the main text. The
structure mirrors the main text: Lemmas A.1--A.4 establish the
conditional distribution machinery; Propositions A.1--A.5 prove
the main results; Proposition A.5 states the linearisation
approximation with explicit error bounds. All symbolic results
have been verified in SymPy (see the companion verification
notebook).

\medskip
\noindent\textbf{Standing assumptions throughout the appendix.}
\begin{enumerate}
  \item[(A1)] \textit{Joint Gaussianity.} $\mu_t \sim \mathcal{N}(0,
    \sigma^2_\mu)$ and $\varepsilon_t \sim \mathcal{N}(0,
    \sigma^2_\varepsilon)$ are mutually independent. The return
    innovation is $\varepsilon^r_t = \frac{1}{2}\mu_t +
    \varepsilon_t$.
  \item[(A2)] \textit{Deviation recursion.} $D_t = D_{t-1}/(1+g) +
    A^r\sigma^2_\mu + B^r\mu_t\varepsilon_t$, where $r \in \{+,-\}$
    denotes the regime determined by $\operatorname{sign}(D_{t-1})$,
    and $A^r$, $B^r$ are the regime-specific coefficients of
    equations~(38)--(39).
  \item[(A3)] \textit{Isserlis.} For any zero-mean bivariate Gaussian
    $(X,Y)$: $\mathbb{E}[X^2Y^2] = \operatorname{Var}(X)
    \operatorname{Var}(Y) + 2\operatorname{Cov}(X,Y)^2$.
\end{enumerate}

% ────────────────────────────────────────────────────────────
\subsection{Conditional distribution lemmas}
% ────────────────────────────────────────────────────────────

\begin{lemma}[Conditional distribution of $(\mu_t, \varepsilon_t)$
  given $\varepsilon^r_t$]
\label{lem:cond_dist}
Under Assumption~(A1), the conditional distribution is bivariate
Gaussian with moments:
\begin{align}
  \mathbb{E}[\mu_t \mid \varepsilon^r_t]
    &= \rho_\mu\,\varepsilon^r_t, \qquad
  \mathbb{E}[\varepsilon_t \mid \varepsilon^r_t]
    = \rho_\varepsilon\,\varepsilon^r_t
    \label{eq:cond_means} \\
  \operatorname{Var}(\mu_t \mid \varepsilon^r_t)
    &= \frac{4\sigma^2_\varepsilon\sigma^2_\mu}
       {4\sigma^2_\varepsilon + \sigma^2_\mu}, \qquad
  \operatorname{Var}(\varepsilon_t \mid \varepsilon^r_t)
    = \frac{\sigma^2_\varepsilon\sigma^2_\mu}
       {4\sigma^2_\varepsilon + \sigma^2_\mu}
    \label{eq:cond_vars} \\
  \operatorname{Cov}(\mu_t, \varepsilon_t \mid \varepsilon^r_t)
    &= \frac{-2\sigma^2_\varepsilon\sigma^2_\mu}
       {4\sigma^2_\varepsilon + \sigma^2_\mu}
    \label{eq:cond_cov}
\end{align}
where $\rho_\mu \equiv \sigma^2_\mu/(2\bar\sigma^2)$,
$\rho_\varepsilon \equiv \sigma^2_\varepsilon/\bar\sigma^2$, and
$\bar\sigma^2 \equiv \sigma^2_\mu/4 + \sigma^2_\varepsilon$.
\end{lemma}

\begin{proof}
By standard Gaussian conditioning. Since $(\mu_t, \varepsilon_t)$
are jointly Gaussian and independent, and $\varepsilon^r_t =
\frac{1}{2}\mu_t + \varepsilon_t$ is a linear combination,
$(\mu_t, \varepsilon_t, \varepsilon^r_t)$ is multivariate Gaussian.
The cross-covariances are $\operatorname{Cov}(\mu_t, \varepsilon^r_t)
= \sigma^2_\mu/2$ and $\operatorname{Cov}(\varepsilon_t,
\varepsilon^r_t) = \sigma^2_\varepsilon$. Dividing by
$\operatorname{Var}(\varepsilon^r_t) = \bar\sigma^2$ gives
\eqref{eq:cond_means}. The conditional variances and covariance
follow from the Schur complement formula applied to the joint
covariance matrix. \textit{Consistency check:}
$\frac{1}{2}\rho_\mu + \rho_\varepsilon =
\sigma^2_\mu/(4\bar\sigma^2) + \sigma^2_\varepsilon/\bar\sigma^2
= \bar\sigma^2/\bar\sigma^2 = 1$, confirming
$\mathbb{E}[\varepsilon^r_t \mid \varepsilon^r_t] = \varepsilon^r_t$.
\qed
\end{proof}

\begin{lemma}[Conditional expectation of $\mu_t\varepsilon_t$]
\label{lem:cond_mueps}
Under Assumption~(A1):
\begin{equation}
  \mathbb{E}[\mu_t\varepsilon_t \mid \varepsilon^r_t]
  = \rho_\mu\bigl[\rho_\varepsilon(\varepsilon^r_t)^2
    - \sigma^2_\varepsilon\bigr]
  \label{eq:cond_mueps}
\end{equation}
This expression is \emph{even} in $\varepsilon^r_t$: it depends only
on $(\varepsilon^r_t)^2$, not on the sign of $\varepsilon^r_t$.
\end{lemma}

\begin{proof}
Decompose $\mu_t\varepsilon_t = \operatorname{Cov}(\mu_t,
\varepsilon_t \mid \varepsilon^r_t) + \mathbb{E}[\mu_t \mid
\varepsilon^r_t]\cdot\mathbb{E}[\varepsilon_t \mid \varepsilon^r_t]$.
Substituting \eqref{eq:cond_cov} and \eqref{eq:cond_means}:
\[
\mathbb{E}[\mu_t\varepsilon_t \mid \varepsilon^r_t]
= -\frac{2\sigma^2_\varepsilon\sigma^2_\mu}
       {4\sigma^2_\varepsilon + \sigma^2_\mu}
  + \rho_\mu\varepsilon^r_t \cdot \rho_\varepsilon\varepsilon^r_t
= -\rho_\mu\sigma^2_\varepsilon
  + \rho_\mu\rho_\varepsilon(\varepsilon^r_t)^2
= \rho_\mu\bigl[\rho_\varepsilon(\varepsilon^r_t)^2
  - \sigma^2_\varepsilon\bigr].
\]
Evenness follows immediately since the right-hand side contains only
$(\varepsilon^r_t)^2$.
\qed
\end{proof}

\begin{lemma}[$\operatorname{Var}(\mu_t\varepsilon_t \mid
  \varepsilon^r_t)$ via Isserlis]
\label{lem:var_mueps}
Under Assumptions~(A1) and (A3):
\begin{equation}
  \operatorname{Var}(\mu_t\varepsilon_t \mid \varepsilon^r_t)
  = \Sigma\bigl[H_0 + H_2(\varepsilon^r_t)^2\bigr]
  \label{eq:var_mueps}
\end{equation}
where $\Sigma \equiv \sigma^2_\varepsilon\sigma^2_\mu /
(4\sigma^2_\varepsilon + \sigma^2_\mu)^2$,
$H_0 \equiv 8\sigma^2_\varepsilon\sigma^2_\mu$, and
$H_2 \equiv 4(16\sigma^4_\varepsilon + \sigma^4_\mu) /
(4\sigma^2_\varepsilon + \sigma^2_\mu)$. In particular
$\operatorname{Var}(\mu_t\varepsilon_t \mid \varepsilon^r_t) > 0$
and both $H_0 > 0$ and $H_2 > 0$.
\end{lemma}

\begin{proof}
Let $u \equiv \mu_t - \mathbb{E}[\mu_t \mid \varepsilon^r_t]$ and
$v \equiv \varepsilon_t - \mathbb{E}[\varepsilon_t \mid
\varepsilon^r_t]$ be the conditional residuals, with
$(u,v) \sim \mathcal{N}(0, V)$ where $V$ is the $2\times2$
conditional covariance matrix of Lemma~\ref{lem:cond_dist}. Denote
conditional means $m_\mu = \rho_\mu\varepsilon^r_t$ and
$m_\varepsilon = \rho_\varepsilon\varepsilon^r_t$. Then
$\mu_t\varepsilon_t = (u+m_\mu)(v+m_\varepsilon)$. Expanding:
\[
\mathbb{E}[(\mu_t\varepsilon_t)^2 \mid \varepsilon^r_t]
= \mathbb{E}[u^2v^2] + m^2_\varepsilon\operatorname{Var}(u)
  + m^2_\mu\operatorname{Var}(v)
  + (m_\mu m_\varepsilon)^2
  + 2m_\mu m_\varepsilon\operatorname{Cov}(u,v),
\]
where all odd moments of $(u,v)$ vanish by Gaussianity. Applying
Assumption~(A3): $\mathbb{E}[u^2v^2] = \operatorname{Var}(u)
\operatorname{Var}(v) + 2\operatorname{Cov}(u,v)^2$. Subtracting
$(\mathbb{E}[\mu_t\varepsilon_t \mid \varepsilon^r_t])^2$ and
collecting terms in $(\varepsilon^r_t)^0$ and $(\varepsilon^r_t)^2$
yields \eqref{eq:var_mueps}. Positivity of $H_0$ and $H_2$ is
immediate from the positivity of all parameters.
\qed
\end{proof}

\begin{lemma}[Slope cancellation]
\label{lem:slope_cancel}
The partial derivative of the standardised threshold $z_t =
-\mu^r_{D_t}/\psi^r_t$ with respect to $(\varepsilon^r_t)^2$ is
regime-invariant:
\[
\frac{\partial z_t}{\partial(\varepsilon^r_t)^2}
= -\frac{\rho_\mu\rho_\varepsilon}{\sqrt{\Sigma H_0}}
  \qquad \text{for both } r \in \{+,-\}.
\]
\end{lemma}

\begin{proof}
From Lemma~\ref{lem:cond_mueps}, $\delta^r \equiv \partial
\mu^r_{D_t}/\partial(\varepsilon^r_t)^2 = B^r\rho_\mu\rho_\varepsilon$.
The conditional standard deviation of $D_t$ at $\varepsilon^r_t = 0$
is $\psi^r_0 = B^r\sqrt{\Sigma H_0}$. Therefore
$\partial z_t/\partial(\varepsilon^r_t)^2 = -\delta^r/\psi^r_0 =
-B^r\rho_\mu\rho_\varepsilon/(B^r\sqrt{\Sigma H_0}) =
-\rho_\mu\rho_\varepsilon/\sqrt{\Sigma H_0}$,
where $B^r$ cancels exactly in both regimes.
\qed
\end{proof}

% ────────────────────────────────────────────────────────────
\subsection{Main propositions}
% ────────────────────────────────────────────────────────────

\begin{proposition}[Asymmetric drift]
\label{prop:asym_drift}
Define the regime-specific stationary points $D^*_r = A^r\sigma^2_\mu(1+g)/g$.
Then $D^*_+/D^*_- = \lambda^2$ and both are strictly positive. The
process $\{D_t\}$ satisfies:
\begin{enumerate}
  \item[(i)] In contraction ($D_{t-1}<0$): $\mathbb{E}[D_t \mid
    D_{t-1}] > D_{t-1}$ unconditionally (submartingale).
  \item[(ii)] In expansion ($D_{t-1}\geq0$): $\mathbb{E}[D_t \mid
    D_{t-1}] \gtrless D_{t-1}$ as $D_{t-1} \lessgtr D^*_+$.
  \item[(iii)] $D^*_+ > D^*_- > 0$: the economy has a positive bias.
  \item[(iv)] The deflation factor $1/(1+g)$ is identical in both
    regimes; asymmetry enters only through the drift $A^r\sigma^2_\mu$.
\end{enumerate}
\end{proposition}

\begin{proof}
By Assumption~(A1), $\mathbb{E}[\mu_t\varepsilon_t] = 0$ (independence
and zero means), so $\mathbb{E}[D_t \mid D_{t-1}] = D_{t-1}/(1+g) +
A^r\sigma^2_\mu$. The fixed point $\mathbb{E}[D_t \mid D_{t-1}] = D_{t-1}$
gives $D^*_r = A^r\sigma^2_\mu(1+g)/g > 0$ since $A^r, \sigma^2_\mu,
g > 0$. The ratio $D^*_+/D^*_- = A^+/A^- = \lambda^2$ follows from
equations~(38)--(39).

\textit{Part (i).} For $D_{t-1} < 0$: $\mathbb{E}[D_t \mid D_{t-1}]
- D_{t-1} = -D_{t-1}g/(1+g) + A^-\sigma^2_\mu > 0$, since both
terms are strictly positive.

\textit{Part (ii).} For $D_{t-1} \geq 0$: the same expression has sign
equal to $\operatorname{sign}(D^*_+ - D_{t-1})$.

\textit{Parts (iii) and (iv).} Both follow directly from the definitions
and $A^+/A^- = \lambda^2 > 1$.
\qed
\end{proof}

\begin{proposition}[Variance asymmetry]
\label{prop:var_asym}\label{prop:Vt_ratio}
Under Assumptions (A1)--(A3), the conditional variance of $D_t$
given $(\varepsilon^r_t, D_{t-1})$ satisfies:
\begin{equation}
  \operatorname{Var}(D_t \mid \varepsilon^r_t, D_{t-1})
  = (B^r)^2\,\Sigma\bigl[H_0 + H_2(\varepsilon^r_t)^2\bigr]
  \label{eq:VarDt}
\end{equation}
and the ratio across regimes is:
\[
\frac{\operatorname{Var}(D_t \mid \varepsilon^r_t, D_{t-1},\,
  D_{t-1}\geq0)}
     {\operatorname{Var}(D_t \mid \varepsilon^r_t, D_{t-1},\,
  D_{t-1}<0)}
= \frac{(B^+)^2}{(B^-)^2} = \lambda^4.
\]
\end{proposition}

\begin{proof}
From Assumption~(A2), $D_t - D_{t-1}/(1+g) - A^r\sigma^2_\mu =
B^r\mu_t\varepsilon_t$. Since $D_{t-1}$ and $\varepsilon^r_t$ are
$\mathcal{F}_{t-1}$-measurable or determined by the current period's
shock distribution, and $B^r$ is $\mathcal{F}_{t-1}$-measurable:
\[
\operatorname{Var}(D_t \mid \varepsilon^r_t, D_{t-1})
= (B^r)^2\operatorname{Var}(\mu_t\varepsilon_t \mid \varepsilon^r_t).
\]
Substituting Lemma~\ref{lem:var_mueps} gives \eqref{eq:VarDt}. The
ratio $(B^+)^2/(B^-)^2 = [\lambda/(2\lambda^0_p)]^2 /
[1/(2\lambda\lambda^0_p)]^2 = \lambda^4$ follows directly.
\qed
\end{proof}

\begin{proposition}[Stationarity condition]
\label{prop:stationarity}
The Threshold-GARCH process for $V_t$ is covariance-stationary if
and only if:
\begin{equation}
  \Phi \;<\; \Phi_{\max} \;\equiv\;
  \frac{2g\lambda^2(g+2)}{(1+g)^2(\lambda^4+1)}
  \label{eq:Phi_max_proof}
\end{equation}
Three corollaries: (a) $\Phi_{\max} \to 0$ as $g \to 0$, recovering
the IGARCH boundary; (b) at $\lambda = 1$, $\Phi_{\max} =
g(g+2)/(1+g)^2 > 0$; (c) $\Phi_{\max}$ is strictly decreasing
in $\lambda$.
\end{proposition}

\begin{proof}
The standard covariance-stationarity condition for a Threshold-GARCH
process is $\alpha + \gamma\pi^* + \beta < 1$, where $\pi^*$ is the
ergodic probability of contraction. At the symmetric stationary point
$\pi^* = 1/2$, this becomes $\alpha + \gamma/2 + \beta < 1$. From
Theorem~1: $\alpha + \gamma/2 = \Phi(\lambda^4+1)/(2\lambda^2)$.
Substituting $\beta = 1/(1+g)^2$ and solving for $\Phi$ gives
\eqref{eq:Phi_max_proof}. Corollary~(a): substitute $g=0$.
Corollary~(b): substitute $\lambda=1$ and simplify. Corollary~(c):
$\partial\Phi_{\max}/\partial\lambda = 2g(g+2)(1+g)^{-2}
\cdot\partial[\lambda^2/(\lambda^4+1)]/\partial\lambda$;
the bracketed expression has derivative
$\lambda(2-2\lambda^4)/(\lambda^4+1)^2 < 0$ for $\lambda > 1$.
\qed
\end{proof}

\begin{proposition}[IGARCH boundary]
\label{prop:igarch}
As $g \to 0$, $\beta = 1/(1+g)^2 \to 1$ and the process converges
to a Threshold-IGARCH with $\alpha + \gamma/2 + \beta = 1$. The
asymmetry coefficient $\gamma$ satisfies $\partial\gamma/\partial g
= 0$: it is independent of the growth rate. A stagnant economy
therefore exhibits the same leverage asymmetry as a growing one,
but with permanently persistent variance shocks.
\end{proposition}

\begin{proof}
$\beta = 1/(1+g)^2 \to 1$ as $g \to 0$ is immediate. That
$\partial\gamma/\partial g = 0$ follows because $\gamma =
\Phi(1-\lambda^4)/\lambda^2$ and $\Phi$ does not depend on $g$
(it contains only $\sigma_\mu$, $\sigma_\varepsilon$, $\lambda^0_p$,
and $\lambda$).
\qed
\end{proof}

% ────────────────────────────────────────────────────────────
\subsection{Linearisation and error bounds}
% ────────────────────────────────────────────────────────────

\begin{proposition}[Linearisation approximation and error bound]
\label{prop:linearisation}
\textbf{Exact part.} For each realised regime $r \in \{+,-\}$,
the conditional variance of $D_t$ is:
\begin{equation}
  V^r_t \equiv \operatorname{Var}(D_t \mid \varepsilon^r_t,
  D_{t-1},\, \text{regime}\,=r)
  = (B^r)^2\,\Sigma\bigl[H_0 + H_2(\varepsilon^r_t)^2\bigr].
  \label{eq:Vt_exact}
\end{equation}
This holds exactly under Assumptions (A1)--(A3).

\medskip
\noindent\textbf{Approximate part.} The expected conditional
variance $\mathbb{E}[V_t \mid \mathcal{F}_{t-1}]$ equals
$\pi_{t-1}V^-_t + (1-\pi_{t-1})V^+_t$, where $\pi_{t-1} =
\Phi(z_{t-1})$ and $z_{t-1}$ is the standardised threshold.
Define $\bar{V} \equiv \Sigma(H_0 + H_2\bar\sigma^2) =
\sigma^2_\varepsilon\sigma^2_\mu$ and expand $\Phi(z_{t-1})$ to
first order around $(D_{t-2} = D^*_r,\, \varepsilon^r_{t-1} = 0)$:
\begin{equation}
  \pi_{t-1} \approx \pi^* +
  \phi(z^*)\left[-\frac{D_{t-2} - D^*_r}{(1+g)\psi^r_0}
  - \frac{\rho_\mu\rho_\varepsilon}{\sqrt{\Sigma H_0}}
    (\varepsilon^r_{t-1})^2\right]
  \label{eq:pi_linearised}
\end{equation}
where $\psi^r_0 = B^r\sqrt{\Sigma H_0}$. Substituting
\eqref{eq:pi_linearised} into $\mathbb{E}[V_t \mid
\mathcal{F}_{t-1}]$ and identifying $V_{t-1} \propto
\pi_{t-2} \cdot [(B^-)^2 - (B^+)^2]\bar{V}$ yields a
recursion of the form:
\begin{equation}
  \mathbb{E}[V_t \mid \mathcal{F}_{t-1}]
  \approx \omega + \alpha(\varepsilon^r_{t-1})^2
  + \gamma(\varepsilon^r_{t-1})^2\mathbf{1}[D_{t-2}<0]
  + \beta V_{t-1}
  \label{eq:Vt_recursion_approx}
\end{equation}
with $\beta = 1/(1+g)^2$ arising from the $D_{t-2}/(1+g)$ term
in $z_{t-1}$, which propagates through two lags to $V_{t-1}$.

\medskip
\noindent\textbf{Error bound.} The error in approximating
$\pi_{t-1}$ by \eqref{eq:pi_linearised} satisfies, by the
Taylor remainder theorem:
\[
|\pi_{t-1} - \pi^{\,\mathrm{approx}}_{t-1}|
\;\leq\;
\frac{\phi(0)}{2\sqrt{e}}\,(z_{t-1} - z^*)^2.
\]
For deviations $|D_{t-2} - D^*_r| \leq \delta\,\psi^r_0(1+g)$,
the implied error in $\mathbb{E}[V_t \mid \mathcal{F}_{t-1}]$ is:
\[
\left|\mathbb{E}[V_t \mid \mathcal{F}_{t-1}]
  - \mathbb{E}[V_t \mid \mathcal{F}_{t-1}]^{\,\mathrm{approx}}\right|
\;\leq\;
\bar{V}\,\bigl|(B^-)^2 - (B^+)^2\bigr|
\,\frac{\phi(0)}{2\sqrt{e}}\,\delta^2.
\]
At $\lambda = 2.25$, $\lambda^0_p = 0.8$, $\sigma_\mu = 1$,
$\sigma_\varepsilon = 0.5$, and $\delta = 0.5$, the relative
error is $5.8\%$. The approximation is accurate when the
economy is within $O(\psi^r_0)$ of the stationary point ---
the empirically relevant regime for moderate deviations from
trend.
\end{proposition}

\begin{proof}
\textit{Exact part.} Follows directly from Proposition
\ref{prop:var_asym} (Lemma~\ref{lem:var_mueps} and
Assumption~(A2)).

\textit{Error bound.} The function $\Phi: \mathbb{R} \to (0,1)$
is twice continuously differentiable with $\Phi'' =
-z\phi(z)$. By Taylor's theorem with remainder:
$|\Phi(z) - \Phi(z^*) - \phi(z^*)(z-z^*)| \leq
\sup_\xi|\Phi''(\xi)| \cdot (z-z^*)^2/2$. Since
$|\Phi''(z)| = |z|\phi(z) \leq \phi(0)/\sqrt{e}$ (attained at
$z = \pm 1$), the bound follows. The remainder maps to the
$V_t$ recursion by multiplying through by $|(B^-)^2 -
(B^+)^2|\cdot\bar{V}$.
\qed
\end{proof}

\subsection{Summary table}

Table~\ref{tab:sympy_verification} records every result
verified symbolically in the companion SymPy notebook
(\texttt{threshold\_garch\_verification\_v3.ipynb}).

\begin{table}[ht]
\centering
\caption{SymPy verification record.}
\label{tab:sympy_verification}
\begin{tabular}{l l l}
\toprule
Block & Result & Status \\
\midrule
1 & Lemma~\ref{lem:cond_dist}: conditional distribution
    $(\mu_t,\varepsilon_t)\mid\varepsilon^r_t$ & PASS \\
2 & Lemma~\ref{lem:cond_mueps}: $\mathbb{E}[\mu_t\varepsilon_t\mid\varepsilon^r_t]$
    even in $\varepsilon^r_t$ & PASS \\
3 & Lemma~\ref{lem:var_mueps}: $\operatorname{Var}(\mu_t\varepsilon_t\mid\varepsilon^r_t)$
    via Isserlis & PASS \\
4 & Proposition~7a: deviation projection $\kappa^r$, $\delta^r$ & PASS \\
5 & Proposition~7: all $\lambda^2$ ratios
    ($\kappa,\,\delta,\,\lambda_p,\,m^*$) & PASS \\
6 & Lemma~\ref{lem:slope_cancel}: slope cancellation & PASS \\
7 & Proposition~\ref{prop:var_asym}: $\operatorname{Var}(D_t)$ ratio $=\lambda^4$ & PASS \\
8 & Theorem~1: exact $\alpha,\gamma,\beta,\omega$ & PASS \\
9 & Proposition~\ref{prop:asym_drift}: asymmetric drift,
    $D^*_+/D^*_- = \lambda^2$ & PASS \\
10 & Proposition~\ref{prop:stationarity}: $\Phi < \Phi_{\max}$ & PASS \\
11 & Proposition~\ref{prop:igarch}: IGARCH boundary,
     $\partial\gamma/\partial g = 0$ & PASS \\
12 & Calibration: minimum $\sigma_\mu/\sigma_u$ at
     Tversky-Kahneman values & PASS \\
\bottomrule
\multicolumn{3}{l}{\footnotesize 51 individual assertions, 0 errors.
  SymPy version 1.14.0.}
\end{tabular}
\end{table}


\section*{Statements and Declarations}

\noindent\textbf{Competing interests.} The author declares no competing interests.

\noindent\textbf{Funding.} No external funding was received for this research.

\noindent\textbf{Data availability.} No datasets were generated or analysed during the current study. The SymPy verification notebook supporting the symbolic results is available from the author on reasonable request.


\bibliographystyle{apalike}
\bibliography{bibliography}

\end{document}
